% Part: lambda-calculus
% Chapter: lambda-definability
% Section: pairs

\documentclass[../../../include/open-logic-section]{subfiles}

\begin{document}

\olfileid{lam}{ldf}{pai}
\olsection{ক্রমযুগল ও পূর্বসূরি}

\begin{defn}
$M$ ও~$N$-এর ক্রমযুগলকে ($\tuple{M,N}$ লেখা হয়) নিচের মতো সংজ্ঞায়িত
করা হয়:
\[
\tuple{M,N} \ident \lambd[f][fMN].
\]
\end{defn}

স্বজ্ঞামূলকভাবে এটি এমন একটি অপেক্ষক, যা একটি অপেক্ষক গ্রহণ করে এবং
সেই অপেক্ষককে ক্রমযুগলটির দুই উপাংশের উপর প্রয়োগ করে। এই ধারণা অনুসারে
আমরা নিচের নির্মাতা পাই; এটি দুটি পদ গ্রহণ করে এবং সেগুলি ধারণকারী
ক্রমযুগলটি ফেরত দেয়:
\[
  \fn{Pair} \ident \lambd[mn][\lambd[f][fmn]]
\]
কোনো ক্রমযুগল দেওয়া থাকলে আমরা তার উপাংশগুলিও ফিরে পেতে চাই। এর জন্য
দুটি অভিক্ষেপ অপেক্ষক দরকার; এগুলি আর্গুমেন্ট হিসেবে একটি ক্রমযুগল গ্রহণ
করে এবং তার প্রথম বা দ্বিতীয় উপাংশ ফেরত দেয়:
\begin{align*}
  \fn{Fst} & \ident \lambd[p][p(\lambd[mn][m])]\\
  \fn{Snd} & \ident \lambd[p][p(\lambd[mn][n])]
\end{align*}

\begin{prob}
  $\fn{Fst}$ ও~$\fn{Snd}$ অভিক্ষেপ অপেক্ষকদুটি কেন কাজ করে, ব্যাখ্যা
  করো।
\end{prob}

ক্রমযুগলের সাহায্যে এবার পূর্বসূরি অপেক্ষককে !!{lambda define} করা যায়:
\[
  \fn{Pred} \ident \lambd[n][\fn{Fst}(n (\lambd[p][\tuple{\fn{Snd}\, {p}, \fn{Succ}(\fn{Snd}\, {p})}]) \tuple{\num 0, \num 0})]
\]
মনে রাখো, $\num n\, f x$ পদটি $f^{n}(x)$-এ হ্রাস পায়; এই ক্ষেত্রে
$f$ এমন একটি অপেক্ষক যা একটি ক্রমযুগল~$p$ গ্রহণ করে এবং $p$-এর দ্বিতীয়
উপাংশ ও সেই দ্বিতীয় উপাংশের উত্তরসূরিকে নিয়ে গঠিত নতুন ক্রমযুগল ফেরত
দেয়; আর $x$ হলো ক্রমযুগল $\tuple{\num{0}, \num{0}}$। ফলে $n=0$ হলে
ফল $\tuple{\num{0}, \num{0}}$, আর অন্যথায়
$\tuple{\num{n-1}, \num n}$। এরপর $\fn{Pred}$ ফলটির প্রথম উপাংশ ফেরত
দেয়।
% BN-SRC-307 TARGET-CORRECTION-BEGIN
\par\smallskip\noindent\textnormal{\small\textbf{উৎস-সংশোধন (BN-SRC-307):} হিমায়িত ব্যাখ্যায় প্রারম্ভিক অবস্থা ও~$n=0$-এর ফলকে দুবার $\tuple{0,0}$ লেখা হয়েছে, যদিও বিশুদ্ধ ল্যাম্বডা ক্যালকুলাসে কাঁচা স্বাভাবিক সংখ্যা পদ নয় এবং ঠিক আগের সংজ্ঞাকারী পদটি $\tuple{\num 0,\num 0}$ ব্যবহার করে। বাংলা ব্যাখ্যায় উভয় উপাংশকে চার্চ সংখ্যাপদ করে $\tuple{\num{0},\num{0}}$ পুনরুদ্ধার করা হয়েছে।} % BN-SRC-307 TARGET-CORRECTION-END

$a$-এর উপর $\fn{Pred}$-কে $b$ বার প্রয়োগ করে ছাঁটা বিয়োগকে
সংজ্ঞায়িত করা যায়:
\[
\fn{Sub} \ident \lambd[ab][b \fn{Pred}\, a].
\]

\end{document}
